% 第3章关键模块实现。
%
% 按WaterOS一级组件的职责边界组织为平台、任务、MM、IPC、VFS、FS、驱动、
% 网络与系统调用九节。每节依次说明职责、接口、数据结构、实现机制和设计取舍，
% 并从当前源码摘录一段能够代表实现方式的代码。

\chapter{关键模块实现}

本章按照\wateros{}的一级组件划分关键实现。平台、任务、内存、IPC、VFS、文件系统、
设备驱动、网络协议栈和系统调用各自拥有明确的状态与接口；跨组件流程只在承担相应职责
的位置展开。表\ref{tab:chapter3-boundary}给出本章九节与源码目录的对应关系。

\begin{longtable}{L{3.2cm}L{4.2cm}L{5.9cm}}
  \caption{关键模块与源码目录}\label{tab:chapter3-boundary}\\
  \toprule
  模块 & 主要目录 & 本章关注内容 \\
  \midrule
  \endfirsthead
  \toprule
  模块 & 主要目录 & 本章关注内容 \\
  \midrule
  \endhead
  平台与架构 & \code{wateros-platform} & 启动、trap、时间、中断、IPI、SMP和架构上下文。 \\
  任务管理 & \code{wateros-task} & 任务、进程线程、调度、等待、退出与回收。 \\
  内存管理 & \code{wateros-mm} & 页帧、地址空间、VMA、缺页、COW和用户内存访问。 \\
  IPC & \code{wateros-ipc} & waitqueue、pipe、futex、signal和共享内存。 \\
  虚拟文件系统 & \code{wateros-vfs} & 路径、挂载、打开文件描述、fd会话和页缓存。 \\
  文件系统 & \code{wateros-fs} & ext4、rootfs、procfs、devfs及文件系统后端。 \\
  设备驱动 & \code{wateros-driver} & 设备发现、注册和块、字符、网卡、输入、显示设备。 \\
  网络协议栈 & \code{wateros-network} & IP配置、TCP/UDP、socket状态和网络事件。 \\
  系统调用 & \code{wateros-syscall}、\code{wateros-abi} & 调用号、参数、分发、用户拷贝和errno。 \\
  \bottomrule
\end{longtable}

以下各节统一采用“职责与边界、接口设计、核心数据结构、实现机制、设计取舍”的顺序。
代码均摘自当前源码，只省略与讨论无关的分支，并用注释标出省略位置。

\section{平台与架构支持}

\subsection{职责与边界}

\code{wateros-platform}负责把固件、指令集和QEMU平台差异收束为上层可调用的统一能力，
包括启动参数、CPU编号、从核启动、核间通知、时钟、中断、页表切换、TLB操作和系统
复位。平台层不维护进程、VMA、文件或socket等内核对象；它只提供操作硬件及架构状态
所需的边界。

RISC-V64路径对接OpenSBI和QEMU virt，LoongArch64路径对接对应的架构入口、CSR和IOCSR
中断机制。早期控制台、panic处理和堆分配器由\code{wateros-runtime}提供，使复杂组件
初始化前已经具备输出错误和分配内存的条件。

\subsection{接口设计}

平台接口进一步区分机器级能力和架构级能力。\code{PlatformBootArgs}描述固件传入的
启动槽位，\code{PlatformSmp}描述会离开当前CPU的SMP操作；\code{ArchTaskContext}、
\code{TrapFrameRead}和\code{TrapFrameWrite}描述任务及trap上下文；\code{ArchTime}与
\code{PlatformTime}提供计时能力。

这些trait按能力而不是按平台型号拆分。\code{PlatformSmp}只包含启动CPU、查询状态、发送
IPI和远端TLB失效等机器级操作；本地中断pending位、trap frame寄存器访问和页表token
切换分别留在架构接口。这样，RISC-V平台实现可以把SMP操作映射到OpenSBI，LoongArch64
实现可以映射到IOCSR或平台中断控制器，而task和MM看到的调用约束保持不变。

平台能力均由编译期选定的零尺寸实现类型提供，trait方法不接收\code{self}。这一形式表达
“当前内核只有一个活动平台”，并避免为每次IPI或TLB操作传递全局对象。可能由固件拒绝的
操作统一返回\code{PlatformSmpResult}；上层据此区分参数错误、能力不支持和固件错误，
而不是依赖某个SBI错误码。

\subsection{核心数据结构}

\code{CpuId}与\code{CpuMask}表示逻辑CPU及CPU集合；\code{HartStatus}表示固件观察到的
处理器状态；\code{IpiKind}记录一次软件层核间请求的原因。固件报告的
\code{Started}不等于内核已经将该CPU标记为online，二者分别表示硬件启动进度与OS初始化
进度。

下面摘录自
\code{wateros-platform/platform-api/api-v0/src/smp.rs}。接口只传递CPU集合和操作结果，
不向上暴露SBI或IOCSR寄存器布局。

\begin{rustcode}
#[repr(u8)]
#[derive(Clone, Copy, Debug, PartialEq, Eq)]
pub enum IpiKind {
    Reschedule = 1 << 0,
    TlbShootdown = 1 << 1,
    TaskNotify = 1 << 2,
}

pub trait PlatformSmp {
    fn start_cpu(
        cpu: CpuId,
        start_addr: usize,
        opaque: usize,
    ) -> PlatformSmpResult<()>;

    fn cpu_status(cpu: CpuId) -> PlatformSmpResult<HartStatus>;
    fn configured_cpu_mask() -> CpuMask;
    fn send_ipi(mask: CpuMask) -> PlatformSmpResult<()>;
    fn flush_tlb_remote(mask: CpuMask) -> PlatformSmpResult<()>;
    fn init_ipi() -> PlatformSmpResult<()>;
}
\end{rustcode}

\subsection{实现机制}

启动CPU完成全局地址空间、堆、任务、驱动和文件系统等共享服务初始化。从CPU使用独立
启动栈进入Rust，依次建立本核trap、timer、IPI接收状态、CPU-local槽位和idle task，
完成后才由task组件发布到online集合。只允许执行一次的初始化由发布状态保护，避免
从CPU重复扫描设备或重复挂载文件系统。

核间请求采用“先写软件状态，再发送硬件通知”的顺序。调度请求在发送IPI前已经写入
目标CPU的pending状态；TLB shootdown在页表更新对其他CPU可见后才发出。目标CPU在IPI、
timer tick及返回用户态边界均会检查相应状态，使通知合并或延迟不会成为永久丢失请求
的原因。

下面摘录自
\code{wateros-platform/platform-impl/impl-qemu-riscv64-opensbi/src/smp.rs}，展示
RISC-V平台如何实现同一SMP契约。\code{start\_cpu}只表示固件接受启动请求，CPU是否真正
可调度仍以task组件发布的online状态为准；IPI与远端TLB失效分别使用SBI IPI和RFENCE。

\begin{rustcode}
impl PlatformSmp for QemuRiscv64OpenSbiSmp {
    fn start_cpu(
        cpu: CpuId,
        start_addr: usize,
        opaque: usize,
    ) -> PlatformSmpResult<()> {
        if !cpu.fits_capacity(MAX_CPUS) {
            return Err(PlatformSmpError::InvalidCpu);
        }
        result(sbi::hart_start(cpu.raw(), start_addr, opaque))
            .map(|_| ())
    }

    fn send_ipi(mask: CpuMask) -> PlatformSmpResult<()> {
        let hart_mask = sbi::HartMask::from_mask_base(
            mask.bits() as usize,
            0,
        );
        result(sbi::send_ipi(hart_mask)).map(|_| ())
    }

    fn flush_tlb_remote(mask: CpuMask) -> PlatformSmpResult<()> {
        let hart_mask = sbi::HartMask::from_mask_base(
            mask.bits() as usize,
            0,
        );
        result(sbi::remote_sfence_vma(
            hart_mask,
            0,
            usize::MAX,
        ))
            .map(|_| ())
    }
}
\end{rustcode}

这里没有把软件pending位写入平台实现：调度原因属于task，待确认的TLB事务属于MM。平台
实现仅负责把已经发布的软件请求送到目标CPU，从而保持trait的硬件边界。

\subsection{设计取舍}

平台实现由Cargo feature在编译期选择，而不是在运行时保留两套架构分支。这减少了裸机
内核中的无效代码路径，也让平台接口成为编译期契约。SMP接口只负责跨CPU操作，本地
调度状态仍由task维护，页表与地址空间生命周期仍由MM维护，从而避免平台层逐渐变成
保存所有全局状态的中心模块。

\designrule{CPU只有在本核trap、timer、IPI和idle task就绪后才能进入online集合；跨核
请求必须先发布软件状态，再发送通知。}

\section{任务管理与调度}

\subsection{职责与边界}

\code{wateros-task}管理可运行任务、用户进程与线程之间的关系，负责创建、入队、切换、
阻塞、唤醒、退出和父进程回收。任务组件维护“谁能够运行、当前在哪个CPU运行、处于
什么生命周期状态”，但不保存地址空间内部页表、打开文件内容或信号队列。

调度器与进程语义位于同一一级组件的不同层次：调度器面向TCB和per-CPU运行队列，进程
注册表面向PID、线程组、父子关系及wait语义。二者通过稳定的任务ID连接。

\subsection{接口设计}

task API以数据类型和函数入口为主，不为了形式统一而设置一个覆盖全部行为的大trait。
\code{TaskId}用于内核调度实体，\code{ProcessId}和\code{ThreadId}用于Linux可见身份；
\code{CloneFlags}描述资源共享方式；\code{TaskSnapshot}与\code{ProcessSnapshot}向
procfs和调试工具提供脱离内部锁的只读状态。

调度策略通过\code{SchedPolicy}、\code{Priority}和\code{Nice}表达。具体策略位于
\code{task-scheduler}子组件，任务生命周期入口不依赖某个特定就绪队列算法。当前
\code{MultiClassScheduler}是调度器内部类型，不通过\code{dyn Trait}暴露；跨组件只调用
\code{spawn}、\code{block\_current}、\code{wake}和\code{schedule\_tick}等语义入口。

这里选择函数式门面有两个原因。其一，调度器是唯一的全局内核服务，不存在运行时切换
多个调度器实例的需求；其二，上下文切换要求调用方同时满足关中断、持有调度临界区等
约束，直接暴露可变调度器对象反而容易绕过这些入口。对外可扩展的部分是策略数据和任务
属性，而不是TCB、运行队列或锁的可变引用。

\subsection{核心数据结构}

\code{TaskControlBlock}保存调度实体和架构上下文，\code{ProcessControlBlock}保存线程组
及进程资源引用，\code{ProcessRegistry}提供PID/TID查询。\code{TaskState}表示调度状态，
\code{ProcessState}则表示进程级生命周期。两者不能合并，因为一个多线程进程可能已经
进入退出阶段，但仍有线程尚未离开CPU。

下面的状态定义摘自
\code{wateros-task/task-api/api-v0/src/task.rs}和\code{process.rs}。

\begin{rustcode}
pub enum TaskState {
    Ready,
    Running,
    Blocking(TaskWaitTarget),
    Sleeping { wake_tick: TaskTick },
    Exited(TaskExitCode),
}

pub enum ProcessTaskState {
    Runnable,
    Exited(TaskExitCode),
}

pub enum ProcessState {
    Running,
    Stopped { signo: u8 },
    Exiting(TaskExitCode),
    Exited(TaskExitCode),
}
\end{rustcode}

\subsection{实现机制}

每个online CPU拥有current task、idle task和调度请求状态。任务进入
\code{Running}前必须从就绪结构移除，并在同一调度临界区记录CPU归属。block、sleep、
exit、yield和时间片耗尽均通过统一状态转换离开运行态；真正切换上下文前释放调度器锁，
避免把旧任务的临界区带到新任务。

就绪队列按调度类别保存任务，per-CPU状态记录当前任务、idle task、是否需要重调度及本地
时间。入队时同时考虑affinity和online CPU集合；如果新任务应在另一CPU运行，调度器先
设置该CPU的重调度标记，再由平台层发送IPI。CPU下线或affinity变化时，任务必须先迁出
不再允许的运行队列，不能只修改一个可见的CPU掩码。

\syscall{fork}和\syscall{clone}先准备地址空间、fd、cwd、凭证和信号对象，全部成功后
才把新任务发布到就绪队列。\syscall{execve}先装载新映像并准备用户栈，再在短临界区
替换当前进程资源；旧地址空间在锁外销毁。\syscall{exit\_group}先把进程发布为
\code{Exiting}，每个线程分别退出，最后一个线程离开后才转为\code{Exited}并允许父进程
回收。

timer tick除了推进调度时间，还会处理延迟或合并的本地重调度请求。下面摘录自
\code{task-scheduler/scheduler-impl/impl-multi-class/src/lib.rs}：

\begin{rustcode}
pub fn schedule_tick() {
    let guard = InterruptGuard::new();
    let cpu_id = cpu::current_cpu_id();
    publish_schedule_reason(cpu_id, 4);
    let (switch_pair, targets) = with_scheduler(|scheduler| {
        let mut switch_pair =
            scheduler.schedule(ScheduleReason::Tick, cpu_id);
        let mut targets = scheduler.take_pending_reschedule_cpus();
        targets.remove(cpu_id);
        let local_requested = scheduler.take_need_resched(cpu_id);
        if switch_pair.is_none() && local_requested {
            switch_pair =
                scheduler.schedule(ScheduleReason::Reschedule, cpu_id);
        }
        (switch_pair, targets)
    });

    dispatch_reschedules(targets, cpu_id);
    if let Some(switch_pair) = switch_pair {
        switch_and_unlock(guard, switch_pair);
    }
}
\end{rustcode}

上段省略了源码注释，保留了完整控制顺序：推进本地调度、取出远端通知目标、消费本地
请求、发送远端IPI，最后按需要切换。

\subsection{设计取舍}

任务状态、CPU归属和就绪结构在同一调度边界内提交，以换取清晰的不变量。跨组件读取
任务状态时使用快照，而不把TCB内部可变引用暴露给procfs或syscall。较重的地址空间销毁、
文件关闭和页回收放在注册表锁外完成，减少长临界区对其他CPU的影响。

\designrule{同一任务在任意时刻只能属于一个CPU或一个非运行状态；退出状态先于唤醒
发布；大型资源销毁不得在调度器或进程注册表锁内执行。}

\section{内存管理}

\subsection{职责与边界}

\code{wateros-mm}管理物理页帧、内核和用户地址空间、VMA、页表、缺页处理、用户内存
访问、ELF装载及用户栈准备。文件映射的数据来源由VFS提供，任务组件只持有地址空间
句柄，平台组件只负责安装页表token和执行TLB操作。

用户地址空间可能被多个CPU上的同进程线程同时访问，也会同时参与系统调用、缺页、
fork和exec。VMA与PTE修改因此由地址空间对象统一同步，不能依赖调用者处于单核环境。

\subsection{接口设计}

\code{AddressSpaceOps}定义页表无关的映射、解除映射、权限修改和地址翻译；
\code{MmapOps}在此基础上增加匿名、文件和设备映射、缺页处理、\syscall{mprotect}、
\syscall{mremap}与\syscall{msync}；\code{UserMemoryOps}提供受检查的用户缓冲区访问。

文件按需映射通过\code{DemandPageLoader}连接VFS与MM。MM只要求加载指定文件偏移的一页，
不依赖具体文件系统或页缓存类型；VFS负责保持文件对象在映射存活期间有效。设备映射则
持有\code{DeviceMappingLease}，其职责不是执行I/O，而是保证用户VMA存在时底层显存或
DMA缓冲区不会被驱动提前释放。

各trait的粒度对应不同的所有权边界。\code{AddressSpaceOps}由task、trap和装载器共同使用；
\code{MmapOps}需要帧分配器并修改VMA，只在MM受控入口中使用；\code{UserMemoryOps}面向
syscall，只暴露经过权限检查的复制操作。\code{DemandPageLoader}使用trait object保存于
VMA中，是因为文件后备对象需要跨越\syscall{mmap}调用的生命周期，并可能在fork时复制。

下面摘录自\code{wateros-mm/mm-api/api-v0/src/mmap.rs}。loader既支持把数据复制到新页，
也允许只读映射直接取得页缓存中的共享物理页；写回与flush只有共享可写映射需要覆盖。

\begin{rustcode}
pub trait DemandPageLoader {
    fn duplicate_box(
        &self,
    ) -> MmResult<Box<dyn DemandPageLoader>>;

    fn mapping_kind(&self) -> DemandMappingKind {
        DemandMappingKind::File
    }

    fn load_page(
        &mut self,
        file_offset: usize,
        dst: &mut [u8],
    ) -> MmResult<()>;

    fn load_shared_page(
        &mut self,
        _file_offset: usize,
    ) -> MmResult<Option<PhysPageNum>> {
        Ok(None)
    }

    fn write_page(
        &mut self,
        _file_offset: usize,
        _src: &[u8],
    ) -> MmResult<()> {
        Err(crate::error::MmError::Unsupported)
    }

    fn flush(&mut self) -> MmResult<()> { Ok(()) }
}
\end{rustcode}

\subsection{核心数据结构}

\code{VirtAddr}、\code{PhysAddr}、\code{VirtPageNum}和\code{PhysPageNum}区分不同地址
类别；\code{PagePerm}表示页权限；\code{MmapRequest}保存地址提示、长度、权限、标志和
后备类型；\code{PageFaultAccess}区分读、写和执行缺页。\code{AddressSpaceId}用于架构
地址空间标识。

下面摘录自\code{wateros-mm/mm-api/api-v0/src/address\_space.rs}和\code{mmap.rs}，
展示上层MM语义如何避免依赖Sv39或LoongArch页表项格式。

\begin{rustcode}
pub trait AddressSpaceOps {
    fn satp_value(&self) -> usize;
    fn map_page_to_ppn(
        &mut self,
        vpn: VirtPageNum,
        ppn: PhysPageNum,
        perm: PagePerm,
    ) -> MmResult<()>;
    fn unmap_page_to_ppn(
        &mut self,
        vpn: VirtPageNum,
    ) -> MmResult<Option<PhysPageNum>>;
    fn protect_page(
        &mut self,
        vpn: VirtPageNum,
        perm: PagePerm,
    ) -> MmResult<()>;
    fn translate_addr(
        &self,
        va: VirtAddr,
    ) -> MmResult<Option<PhysAddr>>;
}

pub enum PageFaultAccess {
    Read,
    Write,
    Execute,
}
\end{rustcode}

\subsection{实现机制}

RISC-V64使用\code{Sv39AddressSpace}，LoongArch64使用
\code{LoongArch64AddressSpace}。两种实现分别完成页表walk和PTE编码，但共享VMA区间、
ELF装载、mmap请求与用户拷贝语义。ELF段和文件mmap可以先登记VMA，首次访问时再按页
分配和装入内容。

fork采用copy-on-write。父子地址空间共享只读映射，可写用户页被标为COW；写缺页发生后
再次确认映射与引用计数，必要时分配新页并复制内容，然后更新PTE。若目标地址空间仍
可能在其他CPU上运行，页表写入后对相应CPU执行TLB shootdown，并等待失效完成后再释放
旧页。

用户拷贝首先验证地址区间和访问权限，再按页访问当前地址空间。路径字符串使用有界NUL
扫描，跨页时重新翻译；页表查询或按需装页失败统一转换为可识别的MM错误，由syscall层
映射为\code{EFAULT}等errno。

缺页入口按VMA种类逐级处理，而不是在trap层判断栈、堆或文件后备。下面摘录自
\code{wateros-mm/mm-impl/impl-sv39/src/user\_heap\_mmap.rs}：

\begin{rustcode}
fn handle_page_fault<A>(
    &mut self,
    allocator: &mut A,
    fault_addr: VirtAddr,
    access: PageFaultAccess,
) -> MmResult<bool>
where
    A: PhysicalFrameAllocator<FrameId = PhysPageNum>,
{
    if self.handle_stack_page_fault(
        allocator,
        fault_addr,
        access,
    )? {
        return Ok(true);
    }
    if self.handle_brk_page_fault(
        allocator,
        fault_addr,
        access,
    )? {
        return Ok(true);
    }
    self.handle_lazy_page_fault(allocator, fault_addr, access)
}
\end{rustcode}

返回值区分“已处理”和“该地址不属于可恢复缺页”，错误则保留分配、装页或页表操作失败的
原因。trap层只负责在处理成功后重试用户指令，未处理的访问再进入用户异常或信号路径。

\subsection{设计取舍}

页表格式放在架构实现中，VMA、COW和用户拷贝留在公共语义层，使双架构能够共享绝大多数
内存管理逻辑。文件映射选择按页加载以避免在\syscall{mmap}中读取整个文件，但VMA因此
必须保存稳定loader和文件偏移。TLB失效按单页、地址空间和跨核目标选择范围，在正确性
边界内减少不必要的全局刷新。

\designrule{页表更新对其他CPU可见后才能发送失效请求；目标CPU完成失效并离开旧执行
上下文前，不得释放相关地址空间或物理页。}

\section{进程间通信与信号}

\subsection{职责与边界}

\code{wateros-ipc}由waitqueue、pipe、futex、signal、共享内存和事件等子组件组成。
IPC负责等待关系、通知状态及IPC对象生命周期，但不自行实现另一套调度器或页表管理。
任务阻塞与唤醒委托task，用户地址及共享页映射委托MM，Linux参数解释委托syscall。

pipe最终表现为VFS文件描述符；futex服务线程同步；signal连接trap返回路径与进程状态；
SysV共享内存则把独立物理页对象映射进多个地址空间。

\subsection{接口设计}

\code{IpcWaitQueueOps}是IPC与任务调度之间的薄接口，提供条件等待、带超时等待、唤醒和
requeue。\code{PipeEndpointOps}描述管道端点，\code{SignalSet}和
\code{SignalAction}描述信号mask与处置，\code{ShmRegistry}管理共享段。

futex队列键不能仅使用用户虚拟地址。private futex还需要包含地址空间作用域，shared
futex则必须使用MM解析出的稳定共享身份，保证不同进程对同一共享字得到相同键。

等待队列trait不暴露任务状态或调度器锁，只接受一个短小的条件闭包。闭包由task实现放入
调度临界区再次执行，从而把“检查条件”和“把当前任务登记为等待者”组成一个不可被wake
插入的状态转换。requeue同样由一个接口完成唤醒与迁移，futex组件不需要先取出任务再
逐个加入另一条队列。

下面摘录自
\code{wateros-ipc/ipc-waitqueue/waitqueue-api/api-v0/src/lib.rs}，保留了futex、pipe和
事件对象共同依赖的几项操作。

\begin{rustcode}
pub trait IpcWaitQueueOps: Sized {
    fn new() -> Self;
    fn id(&self) -> WaitQueueId;
    fn wait_target(&self) -> TaskWaitTarget;

    fn wait_current_while<F>(
        &self,
        condition: F,
    ) -> TaskWaitResult
    where
        F: FnOnce() -> bool;

    fn wake_one(&self) -> Option<TaskId>;
    fn wake_all(&self) -> usize;

    fn requeue_to(
        &self,
        target: Self,
        wake_count: usize,
        requeue_count: usize,
    ) -> usize;
}
\end{rustcode}

\subsection{核心数据结构}

\code{FutexKey}由地址身份、private标记及地址空间作用域组成；
\code{FutexWaitOutcome}表示唤醒、超时或中断结果；\code{PendingSignal}与
\code{SignalDispatch}描述待处理信号和一次投递；\code{ShmSegmentInfo}与
\code{ShmAttachInfo}描述共享段及单个进程的附着关系。

下面摘录自\code{ipc-futex/futex-api/api-v0/src/key.rs}。构造函数把private与shared
两种等价关系写入类型入口，避免等待路径临时拼接不一致的键。

\begin{rustcode}
#[derive(Clone, Copy, Debug, PartialEq, Eq, Hash)]
pub struct FutexKey {
    pub uaddr: usize,
    pub is_private: bool,
    pub private_scope: usize,
}

impl FutexKey {
    pub const fn private(
        uaddr: usize,
        private_scope: usize,
    ) -> Self {
        Self { uaddr, is_private: true, private_scope }
    }

    pub const fn shared(shared_identity: usize) -> Self {
        Self {
            uaddr: shared_identity,
            is_private: false,
            private_scope: 0,
        }
    }
}
\end{rustcode}

\subsection{实现机制}

futex wait先通过MM读取用户值，只有值仍等于期望值时才进入等待队列。条件在调度临界区
内再次检查，以覆盖首次读取与真正阻塞之间的wake竞态。wake、requeue和timeout使用稳定
任务ID操作；线程退出时处理\code{clear\_child\_tid}和robust list，再唤醒用户态等待者。

具体顺序是：由用户地址构造\code{FutexKey}，取得或创建对应队列，读取用户值，随后调用
\code{wait\_current\_while}完成二次检查与阻塞。wake先修改或观察IPC对象状态，再调用
\code{wake\_one}或\code{wake\_all}使任务回到就绪队列；若唤醒发生在二次检查之前，闭包
会发现条件已经变化并取消阻塞，因此不需要依赖IPI恰好到达某个时间点。

signal实现区分进程级pending与线程级pending，并在投递时应用mask和处置。系统调用被
信号中断后，trap返回路径根据调用类型决定返回\code{EINTR}或恢复原调用。signal frame
的寄存器编码由平台\code{SignalFrameCodec}实现，公共信号状态不保存某个架构专用布局。

共享内存注册表拥有段身份和物理页，MM只负责把这些外部所有权页面映射到指定地址空间。
fork增加附件引用，exec与exit解除当前进程映射；标记删除的段在最后一个附件释放后回收。
pipe也采用“预留、复制、提交”的路径：在pipe锁内预留可读区间，释放锁后复制到用户
缓冲区，再根据实际复制长度提交读指针。用户拷贝发生缺页时，只取消或部分提交本次预留，
不会在持有pipe锁时进入MM缺页处理。

\subsection{设计取舍}

waitqueue直接复用任务组件的等待语义，避免IPC内部出现第二套超时和唤醒模型。futex键
显式区分private和shared，使地址相同但地址空间不同的锁不会误相互唤醒。信号投递被放在
trap返回安全点处理，代价是需要保存pending状态，但避免在任意内核临界区直接改写用户
trap frame。

\section{虚拟文件系统}

\subsection{职责与边界}

\code{wateros-vfs}位于系统调用与具体文件系统之间，负责路径规范化、挂载表、命名空间、
打开文件描述、per-task fd会话、工作目录和页缓存。VFS不解析ext4磁盘格式，也不负责
VirtIO块传输；这些能力分别属于FS和driver组件。

普通文件、目录、pipe、控制台、socket和设备节点最终都以VFS句柄进入fd表。对象类型
不同，但dup、fork继承、CLOEXEC、poll和close等文件描述符语义由VFS统一组织。

\subsection{接口设计}

\code{VfsBackend}组合路径打开、挂载、设备清单和根视图能力；\code{VfsIoHandle}定义
已打开对象的读写、seek、flush、ioctl、poll和duplicate入口；\code{VfsOpenOps}把路径
解析结果转换为句柄。\code{VfsFdSession}则面向当前任务管理fd编号及生命周期。

接口使用对象安全的trait object保存不同句柄。只对普通文件有效的操作拥有默认
\code{Unsupported}实现，具体文件、pipe、socket或设备句柄按能力覆盖所需方法。

句柄trait还刻意区分三层状态：fd槽保存\code{CLOEXEC}等描述符标志；打开文件描述保存
dup和fork共享的偏移及状态标志；后端句柄保存文件、pipe或socket对象。\code{duplicate}
复制的是第二层引用，而不是重新按路径open，因此多个fd能够共享文件偏移，两个独立open
则各自拥有偏移。

\subsection{核心数据结构}

\code{NormalizedPath}保存规范化路径；\code{VfsOpenFlags}保存VFS可理解的打开标志；
\code{VfsOpenDescriptionState}表示dup及fork共享的打开文件状态；\code{VfsFd}表示文件
描述符；\code{VfsMetadata}与\code{VfsDirEntry}提供统一元数据和目录项。

下面摘录自\code{wateros-vfs/vfs-api/api-v0/src/backend.rs}和\code{handle.rs}，
展示后端能力组合与统一打开句柄。

\begin{rustcode}
pub trait VfsBackend:
    SingleRootReadView
    + VfsMountOps
    + VfsDevInventory
    + VfsOpenOps
    + VfsMountTable
{
}

pub trait VfsIoHandle: Send + VfsHandleAny {
    fn read(&mut self, buf: &mut [u8]) -> VfsResult<usize> {
        let _ = buf;
        Err(VfsError::Unsupported)
    }
    fn write(&mut self, buf: &[u8]) -> VfsResult<usize> {
        let _ = buf;
        Err(VfsError::Unsupported)
    }
    fn duplicate(&self) -> VfsResult<Box<dyn VfsIoHandle>> {
        Err(VfsError::Unsupported)
    }
    fn open_accmode(&self) -> u32;
    // metadata、seek、flush、ioctl、poll等按句柄能力覆盖。
}
\end{rustcode}

可能阻塞或访问用户内存的读取没有被强行塞入一个\code{read}调用。下面同样摘录自
\code{handle.rs}，\code{prepare\_read}在短临界区内取得可转移的操作对象，
\code{acquire}在fd槽锁外等待或读取数据，最后由\code{finish}依据实际复制进度提交源偏移。

\begin{rustcode}
pub trait VfsReadLease: Send {
    fn bytes(&self) -> &[u8];

    fn finish(
        self: Box<Self>,
        progress: VfsCopyProgress,
    ) -> VfsResult<VfsReadFinish>;
}

pub trait VfsPreparedRead: Send {
    fn acquire(
        self: Box<Self>,
    ) -> VfsResult<Box<dyn VfsReadLease>>;
}

pub trait VfsIoHandle: Send + VfsHandleAny {
    fn prepare_read(
        &mut self,
        _max_len: usize,
    ) -> VfsResult<Box<dyn VfsPreparedRead>> {
        Err(VfsError::Unsupported)
    }
    // 其余方法省略。
}
\end{rustcode}

这一接口不仅服务普通文件，也被pipe、PTY、socket、eventfd、timerfd和inotify复用。各对象
可以有不同的等待条件，但都遵守“准备阶段不长期阻塞、复制失败不错误推进源状态”的契约。

\subsection{实现机制}

路径解析逐段处理cwd、dirfd、符号链接、挂载点和搜索权限，并限制符号链接展开次数。
打开成功后建立稳定的打开文件描述，后续read、write、seek和mmap从句柄操作，不因rename
或unlink重新按原路径查找。每个任务拥有fd会话，fork按语义复制句柄，exec关闭带
\code{CLOEXEC}标志的fd。

路径入口先把相对路径与cwd或dirfd合成为规范化绝对路径，再由最长前缀匹配选择挂载点。
解析目录的过程中逐级检查节点类型和执行权限；最后一段根据open标志决定只查找、创建或
截断。挂载表返回后端及相对路径，VFS桥接层再构造稳定节点句柄。这样，挂载点选择和具体
文件系统的目录查找分别由VFS与FS负责。

页缓存以文件内容身份和页号作为键，协调普通read/write、文件mmap与后端区间I/O。并发
miss只安装一个最终槽位；缓存页取得后还会确认它没有在安装窗口被另一CPU驱逐。脏页
按文件写回，单个文件的后端错误不会污染无关文件的读取结果。

可能阻塞的pipe或poll操作不应长期持有共享fd槽锁。实现提供
\code{with\_current\_io\_detached}，先复制打开文件描述，再在临时句柄上执行阻塞操作；
底层偏移、pipe缓冲和socket状态仍按\code{duplicate}语义共享。

读系统调用使用上述两阶段接口时，先从fd会话取出\code{VfsPreparedRead}，随后释放fd槽
锁并取得lease；数据复制结束后把“已复制字节数”和“是否完成”交给\code{finish}。如果
用户缓冲区在中途发生\code{EFAULT}，普通文件偏移只推进已经复制的部分，pipe和socket也
只消费相同数量的数据。

\subsection{设计取舍}

VFS将“路径对象”和“已打开对象”分开，使目录项变化不会破坏已有fd语义。统一
\code{VfsIoHandle}便于fd表保存多种对象，但能力不匹配必须明确返回
\code{Unsupported}或对应错误，不能静默成功。页缓存位于VFS而非ext4内部，使普通文件
读写和文件mmap能够共享页面身份，同时让具体文件系统只承担区间I/O及元数据操作。

\section{文件系统实现}

\subsection{职责与边界}

\code{wateros-fs}负责具体文件系统后端、根卷选择及内核伪文件系统。当前组件包含FS公共
接口、ext4实现、rootfs管理器、procfs和devfs。FS面向“路径或稳定节点如何映射到文件系统
数据”，VFS面向“用户进程如何通过路径和fd观察这些对象”。

driver组件只注册块设备，FS从devfs取得默认块设备并探测文件系统类型；挂载成功后再由
VFS桥接成统一路径视图。系统调用层不直接引用ext4实现。

\subsection{接口设计}

\code{ReadOnlyFs}和\code{ReadWriteFs}分别描述只读与读写后端，避免只读实现被迫暴露
无意义的写接口。\code{FsImpl}提供名称、探测和挂载入口；\code{RootFsManager}保存当前
根卷；\code{ProcFsView}与\code{DevFsManager}分别提供动态内核视图和设备节点管理。

读写接口同时提供路径操作和稳定节点操作。路径用于创建、查找和目录修改；打开文件后
VFS使用\code{FsNodeId}继续访问，使rename或unlink之后的已有fd仍指向原对象。

\code{FsImpl}和\code{ReadWriteFs}承担不同生命周期：前者是静态登记的实现工厂，只负责
能力查询、卷探测和创建挂载会话；后者是一次挂载产生的有状态对象，保存superblock、
inode缓存及块设备引用。因而多个卷可以由同一个\code{FsImpl}创建多个互不干扰的会话，
VFS也不需要知道当前使用的是哪一种ext4实现。

读写trait以文件系统语义为粒度，不泄露ext4块组或inode内存结构。路径方法便于创建和
目录操作，稳定节点方法保证open后的身份；可选能力使用\code{Unsupported}返回，而不是
由VFS根据文件系统名称进行分支。\code{FsAsyncIo}单独描述可异步提交的I/O，避免普通同步
后端被迫实现无意义的异步状态机。

\subsection{核心数据结构}

\code{FsKind}表示文件系统类型，\code{FsCapability}描述后端能力，
\code{FsNodeType}、\code{FsMetadata}和\code{FsDirEntry}描述节点类型、元数据和目录项，
\code{FsNodeId}保存后端稳定节点身份。\code{SharedRwFs}以共享锁保护可写后端会话。

下面摘录自\code{wateros-fs/fs-api/api-v0/src/handles.rs}。统一注册接口负责能力声明、
块设备探测和挂载，具体文件操作则由返回的\code{SharedFs}或\code{SharedRwFs}完成。

\begin{rustcode}
pub trait FsImpl: Sync {
    fn name(&self) -> &'static str;
    fn supported(&self) -> &'static [FsCapability];

    fn supports(
        &self,
        kind: FsKind,
        mode: FsAccessMode,
    ) -> bool {
        self.supported()
            .iter()
            .any(|c| c.kind == kind && c.access == mode)
    }

    fn probe(
        &self,
        _device: &SharedBlockDevice,
    ) -> FsResult<Option<FsKind>> {
        Ok(None)
    }

    fn mount_ro(
        &self,
        device: SharedBlockDevice,
    ) -> FsResult<SharedFs>;

    fn mount_rw(
        &self,
        _device: SharedBlockDevice,
    ) -> FsResult<SharedRwFs> {
        Err(FsError::Unsupported)
    }
}
\end{rustcode}

\subsection{实现机制}

启动时devfs根据驱动注册表刷新设备节点，FS聚合层选择默认块设备并依次调用已登记实现的
\code{probe}。探测成功后记录活动实现，bring-up阶段再以读写方式挂载rootfs。该顺序把
设备发现、文件系统识别和用户根目录建立拆成可分别诊断的步骤。

ext4实现负责inode、目录项、文件大小、link count和块映射。open建立稳定节点引用，
rename和unlink只改变目录关系；最后一个打开引用关闭后，后端才能回收已经没有目录项的
节点。写入、truncate、mkdir、link和xattr修改均通过统一错误类型返回VFS。

procfs不从磁盘读取固定文本，而是通过注册回调取得任务、挂载、内存、IPC和网络快照，
按访问时状态生成内容。devfs将块、字符、输入和显示设备映射为节点，并把打开操作转换为
对应设备句柄。

下面摘录自\code{wateros-fs/src/lib.rs}。聚合层只根据trait能力和探测结果选择实现，
既不读取ext4 magic，也不直接构造具体后端类型；读写能力优先，只有后端声明不支持时才
接受只读会话。

\begin{rustcode}
let mut chosen: Option<(
    &'static dyn api_v0::FsImpl,
    api_v0::FsKind,
)> = None;

for imp in registered_fs_impls() {
    match imp.probe(&device) {
        Ok(Some(kind))
            if imp.supports(
                kind,
                api_v0::FsAccessMode::ReadWrite,
            ) =>
        {
            chosen = Some((*imp, kind));
            break;
        }
        Ok(Some(kind))
            if imp.supports(
                kind,
                api_v0::FsAccessMode::ReadOnly,
            ) =>
        {
            chosen = Some((*imp, kind));
            break;
        }
        Ok(_) => {}
        Err(err) => logging::warn!(
            "[fs] probe via {} failed: {:?}",
            imp.name(),
            err,
        ),
    }
}
\end{rustcode}

探测错误只淘汰当前实现，不能破坏块设备或阻止后续实现继续识别。同一实现报告的
\code{FsKind}还要与静态能力表相符，避免“能够识别但不能按请求模式挂载”的后端被误选。

\subsection{设计取舍}

FS和VFS分层使磁盘格式实现不承担per-task fd语义，也使procfs、devfs等伪文件系统能够
接入同一命名空间。稳定节点接口增加了后端引用管理，但能够正确表达“目录项已删除而
打开文件仍可访问”的Linux语义。同步写路径明确区分页缓存写回、文件系统同步和块设备
flush，错误沿调用链返回用户态。

\section{设备管理与驱动}

\subsection{职责与边界}

\code{wateros-driver}负责机器级设备发现、设备类型匹配、具体驱动初始化和设备对象注册。
其子组件覆盖块设备、字符设备、网络设备、输入设备和显示设备。驱动层处理寄存器、VirtIO
队列、DMA和中断，不解释ext4、TCP或用户态fd语义。

RISC-V64平台主要从DTB发现VirtIO-MMIO设备；LoongArch64平台通过PCIe配置空间发现
VirtIO-PCI设备。两条路径最终都登记为公共设备trait object。

\subsection{接口设计}

\code{MachineDriver}定义所选机器在内核基础设施就绪后的设备初始化入口；
\code{BlockDevice}、\code{CharacterDevice}、\code{NetworkDevice}、
\code{InputDevice}和\code{DisplayDevice}分别定义设备类别接口。聚合层的
\code{machine()}根据feature返回当前机器实现。

各设备子系统通过\code{supported\_devices()}声明可绑定的\code{compatible}，机器驱动
负责扫描并把匹配节点交给具体实现。上层取得的是共享设备对象，不直接持有平台探测结构。

机器trait与设备trait分开，是因为二者的实例数量和调用阶段不同。\code{MachineDriver}
在启动阶段调用一次，拥有总线枚举和资源解释能力；\code{BlockDevice}等trait则对应每个
已发现设备，在系统运行期间反复被FS、VFS或network调用。注册表保存
\code{Arc<Mutex<Box<dyn Trait>>>}，\code{Arc}管理跨组件生命周期，\code{Mutex}串行化设备
内部队列，trait object负责消除具体VirtIO传输类型。

\subsection{核心数据结构}

\code{DeviceType}区分Block、Character、Network、Display和Input；
\code{DeviceInfo}保存节点名、compatible列表、设备类型、MMIO范围和中断线；
\code{MmioRegion}与\code{IrqLine}描述硬件资源。\code{SharedBlockDevice}等类型使用
共享锁保护动态分派的设备对象。

下面摘录自\code{wateros-driver/driver-api/api-v0/src/lib.rs}，可以看到机器级接口与
单个设备信息之间的边界。

\begin{rustcode}
pub enum DeviceType {
    Block,
    Character,
    Network,
    Display,
    Input,
    Unknown,
}

pub struct DeviceInfo {
    pub node_name: String,
    pub compatible: String,
    pub compatibles: Vec<String>,
    pub device_type: DeviceType,
    pub mmio: Option<MmioRegion>,
    pub irq: Option<IrqLine>,
}

pub trait MachineDriver {
    fn init_after_boot(&self) -> DriverResult<()>;
    fn realtime_ns(&self) -> DriverResult<Option<u64>> {
        Ok(None)
    }
    fn test(&self);
}
\end{rustcode}

设备类别trait只包含该类设备能够可靠提供的语义。下面摘录自
\code{wateros-driver/driver-block/block-api/api-v0/src/lib.rs}；整块读写是必需能力，
任意字节读取由公共默认实现完成，\code{flush}则明确表示此前成功写入必须提交到稳定介质。

\begin{rustcode}
pub trait BlockDevice: Send {
    fn block_size(&self) -> usize { BLOCK_SIZE }
    fn total_blocks(&self) -> Option<u64> { None }

    fn read_blocks(
        &mut self,
        start_block: Lba,
        buf: &mut [u8],
    ) -> DriverResult<()>;

    fn write_blocks(
        &mut self,
        start_block: Lba,
        buf: &[u8],
    ) -> DriverResult<()>;

    fn flush(&mut self) -> DriverResult<()>;
}
\end{rustcode}

块设备接口没有加入路径、页缓存或文件同步概念；相应地，\code{NetworkDevice}只处理完整
二层帧和链路属性，不出现TCP端点。类别trait保持这一边界后，VirtIO-MMIO、VirtIO-PCI或
缓存包装器都可以替换，而上层协议不需要修改。

\subsection{实现机制}

聚合层先选择当前\code{MachineDriver}，随后收集各设备类别的支持条目。RISC-V64实现解析
DTB的\code{compatible}、\code{reg}与中断属性，对VirtIO-MMIO设备读取设备ID；LoongArch64
实现枚举PCI函数并解析VirtIO capability。设备握手成功后，具体实现被包装成对应trait
object并加入全局注册表。

每个匹配项先验证MMIO或PCI资源，再完成VirtIO状态协商、feature确认和队列建立，最后才
注册共享设备对象。上层因此不会观察到只完成一半初始化的设备。注册表下标保持稳定，
devfs刷新时根据类别和下标生成节点；刷新只建立视图，不重新初始化硬件。

块设备通过\code{BlockDevice}提供整块读写和flush，可选缓存层在不改变上层接口的情况下
包装真实设备；字符设备包括UART、RTC及空设备；输入和显示设备通过devfs向用户态暴露。
网卡驱动只提供帧收发能力，协议解析和socket状态由下一节的network组件负责。

初始化失败会记录设备及错误信息，并允许其他设备继续探测。根文件系统是否能够挂载由FS
阶段单独判断，不把“某个非必需设备缺失”直接等同于平台启动失败。

中断路径只确认设备状态并推进相应队列，耗时的数据处理由上层在受控入口继续完成。DMA
缓冲区由具体驱动持有，并在设备对象仍被注册或被VMA lease引用时保持有效；设备对象的
共享生命周期因此先于devfs句柄和用户映射结束。

\subsection{设计取舍}

按设备类别拆分trait，使块设备的flush语义、字符设备的流式I/O和网卡的帧队列不会被塞入
一个巨大的通用设备接口。机器驱动集中处理总线差异，设备类别组件集中处理设备行为。
共享trait object带来一次动态分派，但换取了FS、VFS和network对MMIO/PCI实现的解耦。

\section{网络协议栈}

\subsection{职责与边界}

\code{wateros-network}负责网络配置、协议栈推进、TCP/UDP/ICMP socket、连接状态、收发
缓冲和网络事件。它从driver组件取得\code{NetworkDevice}收发以太网帧，但不负责PCI或
VirtIO队列初始化；它通过VFS把socket包装成fd，但不管理通用fd编号。

因此，网卡属于设备驱动，IP、TCP、UDP和socket属于网络协议栈。两者在
\code{NetworkDevice}边界处连接，章节上也分别说明。

\subsection{接口设计}

network API定义后端无关的地址、端点、socket类型、状态、错误和快照；当前具体协议栈
由\code{impl-smoltcp}提供。聚合层只选择性转发创建socket、bind、connect、listen、
accept、send、receive、poll和选项操作，不把smoltcp内部socket set直接暴露给syscall。

\code{SocketRef}是syscall和VFS使用的主要对象接口。它把底层协议栈句柄放入共享对象，
使dup、fork和正在执行的系统调用可以共同持有同一打开socket。

网络层没有再定义一个复制smoltcp全部能力的“协议栈trait”。当前后端由feature静态选择，
对外变化点是稳定的\code{SocketRef}方法和数据类型；真正需要动态替换的网卡使用上一节的
\code{NetworkDevice} trait。这样可以避免内部TCP状态、socket set借用和设备token穿过
组件边界，同时仍允许syscall层独立实现Linux地址族、标志位和errno转换。

接口中的快照与租约分别解决只读观察和有副作用读取。\code{SocketPollSnapshot}在一次
协议栈临界区内复制可读、可写、连接错误和accept状态，调用方不能持有内部socket引用；
\code{SocketReceiveLease}则持有一次接收预留，只有显式\code{finish}才消费数据。二者都
允许任务在等待期间完全释放协议栈锁。

\subsection{核心数据结构}

\code{NetworkAddress}支持IPv4与IPv6，\code{NetworkEndpoint}保存地址、端口和scope id；
\code{SocketDomain}和\code{SocketKind}表示地址族与类型；\code{SocketState}描述内核侧
状态机；\code{SocketPollSnapshot}在一次协议栈临界区中取得poll所需的完整状态。

下面摘录自\code{wateros-network/src/socket/object.rs}。共享对象的最后一个引用释放时
才关闭底层socket，accept则在同一监听句柄锁内完成连接取出与替换。

\begin{rustcode}
struct SocketShared {
    handle: Mutex<StackSocketHandle>,
    domain: SocketDomain,
    status_flags: AtomicUsize,
}

impl Drop for SocketShared {
    fn drop(&mut self) {
        let handle = *self.handle.get_mut();
        if let Err(err) = stack::socket_close(handle) {
            log::warn!(
                "[socket-ref] final close failed handle={:?} err={:?}",
                handle,
                err,
            );
        }
    }
}

#[derive(Clone)]
pub struct SocketRef {
    inner: Arc<SocketShared>,
    inode: u64,
}

impl SocketRef {
    pub fn accept(
        &self,
        status_flags: usize,
    ) -> NetworkResult<(SocketRef, NetworkEndpoint)> {
        let mut listener = self.inner.handle.lock();
        let (established, replacement, peer) =
            stack::socket_accept(*listener)?;
        *listener = replacement;
        Ok((
            Self::from_stack_handle(
                established,
                self.domain(),
                status_flags,
            ),
            peer,
        ))
    }
}
\end{rustcode}

\subsection{实现机制}

协议栈初始化时取得网卡设备并安装IP地址、前缀和网关。受控poller从设备队列收取帧、
推进协议状态并发送待输出帧；等待网络事件前释放协议栈全局锁，避免任务睡眠期间阻塞
其他socket。

listen和accept维护监听句柄与已建立连接，send根据当前发送容量推进数据，receive使用
预留与提交两阶段接口，避免把用户拷贝放在协议栈锁内。poll/epoll从
\code{SocketPollSnapshot}取得可读、可写、连接完成、连接错误和pending accept状态，
等待后重新获取快照，而不长期保存内部可变引用。

socket作为VFS句柄进入fd表，网络模块仍拥有底层状态和最终close语义。procfs使用
\code{NetworkSocketSnapshot}生成\code{/proc/net}内容，格式化过程只读取快照。

下面摘录自\code{wateros-network/src/socket/lease.rs}。接收前先用快照快速判断状态，
随后在协议栈中预留数据并复制到内核暂存区；lease同时持有\code{SocketRef}，保证用户
复制期间底层socket不会被最后一个fd关闭。

\begin{rustcode}
impl SocketRef {
    pub fn prepare_receive(
        &self,
        max_len: usize,
    ) -> Result<SocketReceiveLease, SocketRecvError> {
        let snapshot = self.poll_snapshot()
            .map_err(|_| SocketRecvError::InvalidSocket)?;
        if !snapshot.can_recv {
            if snapshot.kind == SocketKind::Tcp && !snapshot.may_recv {
                return Err(SocketRecvError::Finished);
            }
            return Err(SocketRecvError::Empty);
        }

        let mut data = Vec::new();
        data.try_reserve_exact(max_len)
            .map_err(|_| SocketRecvError::NoMemory)?;
        data.resize(max_len, 0);
        let reservation = self.with_handle(|handle| {
            stack::socket_prepare_recv(handle, &mut data)
        })?;
        data.truncate(reservation.staged_len());
        Ok(SocketReceiveLease {
            _socket: self.clone(),
            reservation: Some(reservation),
            data,
        })
    }
}
\end{rustcode}

正常路径按实际复制长度提交；若用户拷贝失败或lease提前析构，drop路径以零字节、未完成
状态取消预留。TCP字节流可提交部分数据，UDP数据报则由完成标志决定整报文消费语义。

\subsection{设计取舍}

网络API不直接暴露smoltcp类型，使协议栈实现与Linux socket ABI之间保留转换层。
\code{SocketRef}以共享生命周期表达dup和fork，比按fd复制底层连接更符合打开对象语义。
poll使用一次性快照，代价是唤醒后需要重新检查，但可以避免跨睡眠持有协议栈锁。

\section{系统调用与Linux ABI}

\subsection{职责与边界}

\code{wateros-syscall}是用户态进入内核服务的统一入口，负责识别调用号、读取寄存器参数、
安全访问用户内存、检查标志和权限、调用状态所属组件，并把内部错误转换成Linux errno。
\code{wateros-abi}维护Linux generic 64位调用号及ABI类型，两者分别承担“编号与布局”和
“具体行为”。

系统调用组件不保存进程、文件、地址空间或socket的权威状态。新增调用如果需要新的内核
机制，应先在task、MM、VFS、IPC或network中建立通用接口，再由syscall完成ABI适配。

\subsection{接口设计}

\code{SyscallNumber}封装调用号，\code{SyscallArgs}按ABI顺序保存参数寄存器，
\code{SyscallPacket}组合一次完整请求。\code{UserRet}统一编码非负成功值与负errno，
\code{KernelResult<T>}用于实现内部传递未编码的错误。

trap路径只构造参数并调用\code{dispatch\_syscall\_from\_trap}。具体分发按文件、进程、MM、
IPC、网络、时间和凭证等功能文件组织；是否能够自动重启由单独表记录，避免所有被信号
打断的调用一律重试。

这里同样没有设置覆盖全部系统调用的trait。调用号集合是编译期确定的Linux ABI，使用
直接分发表可以在入口处看到“编号到实现”的一一对应；各\code{sys\_*}函数再依赖所属
组件的稳定接口。可替换实现位于MM、FS、driver等组件内部，syscall层自身不是运行时插件。

每个实现按“ABI解析、权限与范围检查、取得内核对象、执行组件操作、写回用户结果”的顺序
组织。用户指针不会转换为可长期保存的Rust引用，可能睡眠的路径也不会跨等待保留用户页
或fd槽的借用。内部trait返回领域错误，syscall文件在离ABI最近的位置完成errno映射，
使组件接口不依赖Linux错误号。

\subsection{核心数据结构}

\code{ErrNo}封装Linux错误号；stat、statx、socket地址、timespec和signal frame等用户
结构体在ABI或syscall实现层定义，不与内核内部对象共用布局。用户指针始终以整数进入
分发层，再通过MM用户拷贝接口读写。

下面摘录自\code{wateros-syscall/syscall-api/api-v0/src/args.rs}和
\code{return\_value.rs}。数据结构只承载ABI，不解释某个参数槽位的具体含义。

\begin{rustcode}
#[repr(C)]
#[derive(Clone, Copy, Debug)]
pub struct SyscallArgs {
    pub args: [usize; MAX_SYSCALL_ARGS],
}

impl SyscallArgs {
    pub const fn from_regs(
        regs: [usize; MAX_SYSCALL_ARGS],
    ) -> Self {
        Self { args: regs }
    }

    pub const fn arg(&self, idx: usize) -> usize {
        self.args[idx]
    }
}

#[repr(transparent)]
pub struct UserRet(pub isize);

impl UserRet {
    pub const fn from_success(v: usize) -> Self {
        Self(v as isize)
    }
    pub const fn from_error(errno: ErrNo) -> Self {
        Self(errno.user_ret())
    }
}
\end{rustcode}

\subsection{实现机制}

用户态\code{ecall}进入架构trap后，trap frame提供调用号和最多六个参数。分发入口记录
调用统计，按调用号选择实现，并在调用完成且不再持有VFS或设备锁后更新进程I/O统计。
普通调用返回时推进用户PC并写回\code{UserRet}；\syscall{execve}成功后已经换成新映像
的trap frame，不能再按普通返回路径覆盖。

具体实现先检查指针、长度、fd和标志，再调用对应组件。路径和结构体通过有界user copy
进入内核；产生副作用的调用在修改状态前完成权限检查。内部的\code{VfsError}、
\code{MmError}、\code{ProcessError}和\code{NetworkError}分别映射为稳定errno。不支持
的调用或功能返回\code{ENOSYS}，而不是成功占位或直接panic。

procfs相关系统调用使用task、MM、VFS、IPC和network快照生成内容。CPU online信息与
affinity使用同一CPU集合，计时接口使用平台单调时间，避免不同用户接口对同一内核状态
给出互相矛盾的结果。

下面摘录自\code{wateros-syscall/syscall-impl/impl-kernel/src/lib.rs}。统一入口完成分发，
并在调用结束、不再持有VFS或设备锁后记录I/O统计。

\begin{rustcode}
pub fn dispatch_syscall_from_trap(
    syscall_nr: usize,
    syscall_args: SyscallArgs,
) -> isize {
    sys::record_syscall();
    let result = syscall_nr_dispatch::dispatch_syscall_by_nr(
        syscall_nr,
        syscall_args,
    );
    if result >= 0 {
        account_process_io(syscall_nr, result as u64);
    }
    result
}
\end{rustcode}

trap返回部分摘录自\code{os/src/trap\_handler.rs}。这里必须单独处理\syscall{execve}和
可重启调用，不能机械地推进旧trap frame。

\begin{rustcode}
let exec_succeeded = syscall_nr == EXEC && syscall_ret >= 0;
if !exec_succeeded {
    cx.add_user_pc(SYSCALL_INSN_BYTES);
    cx.set_syscall_ret(UserRet(syscall_ret));
    if syscall_ret == syscall::ErrNo::EINTR.user_ret()
        && syscall::is_restartable_syscall(syscall_nr)
    {
        restart = Some((syscall_nr, syscall_args));
    }
}
\end{rustcode}

\syscall{execve}成功时，新映像的入口、栈和寄存器已经写入trap frame，继续使用旧frame的
普通返回逻辑会破坏新上下文。可重启调用则保存原调用号与参数，待信号投递完成后由返回
路径恢复；不在重启表中的调用保持\code{EINTR}，由用户态决定是否重试。

\subsection{设计取舍}

调用号表和行为实现分离，使ABI编号能够保持稳定，同时允许内核机制逐项接入。系统调用
层集中做用户拷贝与errno转换，可以防止相同Linux结构体解析散落到各组件。代价是分发层
需要依赖多个组件，因此通过“只转换、不持有状态”的边界控制其复杂度。

\designrule{成功路径和错误路径同属ABI；参数无效、权限不足或用户地址不可访问时，必须
在产生不可回滚副作用之前返回相应errno。}
